% Algebraic Data Types Are Algebraic
% K. Siddiqui and Claude Sonnet
% LaTeX --- pearl preprint

\documentclass[11pt, a4paper]{article}

% --- Packages ---
\usepackage{amsmath, amssymb, amsthm}
\usepackage{geometry}
\usepackage{microtype}
\usepackage{marginnote}
\usepackage{hyperref}

% --- Page geometry ---
\geometry{
  a4paper,
  top=2.5cm, bottom=2.8cm,
  inner=2.8cm, outer=4.8cm,
  marginparwidth=2.8cm,
  marginparsep=0.5cm
}

% --- Theorem environments ---
\newtheoremstyle{pearlstyle}
  {6pt}{6pt}{\itshape}{0pt}{\bfseries}{.}{0.5em}{}
\theoremstyle{pearlstyle}
\newtheorem*{conjecture}{Conjecture}
\newtheorem*{remark}{Remark}

% --- Notation macros ---
\newcommand{\bone}{\mathbb{1}}
\newcommand{\btwo}{\mathbb{2}}
\newcommand{\bzero}{\mathbb{0}}
\newcommand{\bF}{\mathbb{F}}
\newcommand{\Q}{\mathbb{Q}}
\newcommand{\Z}{\mathbb{Z}}
\newcommand{\N}{\mathbb{N}}
\newcommand{\C}{\mathbb{C}}
\newcommand{\R}{\mathbb{R}}
\newcommand{\Exp}{\mathtt{exp}}
\newcommand{\Sn}{S_n}
\newcommand{\coloneq}{\mathrel{:\!=}}

% --- Spacing ---
\setlength{\parskip}{4pt plus 1pt}
\setlength{\parindent}{18pt}

% ======================================================
\begin{document}

\title{Algebraic Data Types Are Algebraic\\[4pt]
  {\large --- and Transcendental Types Are Transcendental}}
\author{Khalid Siddiqui \and Claude Sonnet}
\date{\itshape Pearl preprint --- \today}
\maketitle

\begin{abstract}
The word ``algebraic'' appears in two apparently unrelated contexts:
algebraic data types in programming language theory, and algebraic numbers in
arithmetic. We show these are not merely homonyms. The type-theoretic operations
of sum, product, and exponentiation generate exactly the class of algebraic
functions, and the values outside that closure --- those unreachable by any
combination of these constructors --- are transcendental in both senses
simultaneously. As a worked illustration, the exponential type former
$\mathtt{exp}$, uniquely characterised as the fixed point of differentiation,
is read as a power series and evaluated along an imaginary axis via
Cayley--Dickson doubling. Its period is inexpressible in the language used to
define it. This is transcendence as a syntactic phenomenon: a mild and
pronounceable form of G\"{o}del incompleteness, visible at the level of types.
\end{abstract}

% ======================================================
\section{A thing being is any thing}

We begin with notation chosen to keep the levels honest. Write $\bone$ for the
unit type --- the type with exactly one inhabitant --- and $\btwo$ for the type
with exactly two. A central confusion in ordinary notation is that ``1'' and
``2'' carry arithmetic connotations (integers, field elements, cardinals) that
obscure what is structurally primitive. Blackboard bold signals that we mean the
types themselves, not numbers that happen to index them.

\medskip
\noindent\textbf{A note on notation.} This paper uses typography to enforce a
distinction that prose tends to elide. Blackboard bold ($\bone$, $\btwo$,
$\bzero$) denotes \emph{types as types} --- structural objects, not numbers that
happen to share a name. Typewriter font ($\Exp$) denotes \emph{type
constructors} --- terms in a formal language, defined by their structural rules,
not by analytic properties. Roman font is reserved for evaluated,
interpreted, model-specific versions. The typographic discipline is the argument
in miniature: the distinction between $\Exp$ and $e^x$ is exactly the
distinction this paper is about. There is only one $\mathtt{exp}$onential.

\medskip
Now consider the quotient $X/X$ for a type $X$. Whatever $X$ is, the type of
``an $X$ without any information about which $X$'' has exactly one inhabitant:
the fact of there being an $X$ at all, without specification. So
\[
  X/X \;=\; \bone.
\]
This is not arithmetic; it is the observation that self-quotients collapse to
the unit. A thing being is any thing.

From this, $\btwo \coloneq \bone + \bone$ --- two distinct units, a binary
choice, a coproduct. No integers have been assumed. The arithmetic will be
derived, not imported.

% ======================================================
\section{The Taylor series is a type}

We define $\Exp$ structurally, without importing any analytic machinery. The
right starting point is not the power series but the property that characterises
$\Exp$ uniquely: it is the fixed point of differentiation.

In the setting of Joyal's combinatorial species, differentiation has a clean
combinatorial meaning. The \emph{derivative} $F'$ of a species $F$ is the
species of $F$-structures with one element \emph{pointed} --- singled out,
marked, distinguished from the rest. This is the ``one-hole context''
construction: an $F'$-structure is an $F$-structure with a distinguished site of
attention.\marginnote{\footnotesize McBride's work on the derivative of a
container type realises exactly this construction in dependent type theory: the
type of one-hole contexts of a container $F$ is its derivative $F'$. The
present observation is that $\Exp$ is the unique fixed point of that operation.}

Now consider the species $\mathbf{E}$ of finite sets: an $\mathbf{E}$-structure
on a set $S$ is simply $S$ itself, unadorned. Pointing one element of a finite
set gives a finite set with a distinguished element --- which is again a finite
set (the distinguished element carrying no extra structure beyond its
singling-out). So $\mathbf{E}' = \mathbf{E}$: the species of finite sets is its
own derivative. It is the fixed point of the pointing operator.

This is the structural definition of $\Exp$:
\[
  D(\Exp) \;=\; \Exp, \qquad \Exp(\bzero) \;=\; \bone.
\]
The initial condition $\Exp(\bzero) = \bone$ --- the empty collection is the
unit --- pins down the unique fixed point. Everything else follows.

\medskip\noindent\textbf{The power series as consequence.} Expanding the fixed
point equation around $\bzero$ forces the coefficients one by one. If
$\Exp(X) = \sum_{n \geq 0} a_n X^n$, then $D(\Exp) = \Exp$ gives
$a_n = a_{n-1}/n$ for all $n \geq 1$, and $a_0 = \bone$ gives $a_n = \bone/n!$.
So
\[
  \Exp(X) \;=\; \sum_{n \geq 0} \frac{X^n}{n!}.
\]
The $n!$ denominators are not analytic conveniences; they are forced by the
fixed point equation together with the initial condition. Their combinatorial
meaning is immediate: $X^n/n!$ is the type of \emph{unordered} $n$-tuples, the
quotient of the ordered product $X^n$ by the action of the symmetric group
$\Sn$, whose groupoid cardinality is exactly
$n!$.\marginnote{\footnotesize Groupoid cardinality: as a groupoid,
$|\Sn| = n!/n! = 1$, so the quotient $X^n/\Sn$ has cardinality $|X|^n/n!$
--- the $n$-th term of an exponential generating function.}

So $X^2/2! = \{X,X\}$ is the unordered pair, and the full series
\[
  \Exp(X) \;=\; \bone \;+\; X \;+\; \{X,X\} \;+\; \{X,X,X\} \;+\; \cdots
\]
is the type of \emph{finite unordered collections} of $X$ --- the free
commutative monoid on $X$, the species $\mathbf{E}$ of finite sets. This is not
an analogy: the ``exponential'' in ``exponential generating function'' is $\Exp$,
and the $n!$ denominators arise \emph{because} the structures are unordered.
Type theorists have been doing transcendence theory all along, and $\Exp$ was
always the fixed point of their own derivative.

% ======================================================
\section{The algebraic and the transcendental}

The type constructors $+$, $\times$, and $(-)^{(-)}$ --- sum, product,
exponentiation --- generate, by composition and fixed point, the class of
\emph{algebraic} types: those whose generating functions satisfy polynomial
equations over $\Q$. A list type $L(X) = \bone + X \times L(X)$ satisfies
$L = 1/(1-X)$, a rational function; a binary tree type satisfies a
quadratic.\marginnote{\footnotesize $L(X) = 1/(1-X)$, convergent for $|X|<1$.
The fixed-point equation $L = \bone + X \cdot L$ gives $L(1-X) = \bone$,
recovering the same rational.} The values of these generating functions at
rational arguments are, in general, algebraic numbers.

The exponential type $\Exp(X)$ is different. Its generating function $e^x$ does
not satisfy any polynomial equation over $\Q$ --- it is transcendental as a
function. And its values at nonzero algebraic arguments are transcendental
numbers, by Lindemann--Weierstrass. The exponential type constructor generates
an object that escapes the algebraic closure of the language used to build it.

This is the first sense in which algebraic data types are algebraic: the
boundary of the algebraically-generated types is precisely the boundary of
algebraic functions, and what lies beyond is transcendental in the
number-theoretic sense.

% ======================================================
\section{Cayley--Dickson and the imaginary axis}

We have been evaluating $\Exp$ at a real variable $X$. To reach $\pi$, we must
evaluate it on the imaginary axis, which requires extending our number system.

The complex numbers $\C$ arise from $\R$ by Cayley--Dickson doubling: one
adjoins a new element with square $-\bone$, and defines multiplication
accordingly on pairs. This is a \emph{structural} construction --- it requires
no geometric intuition about rotation or the complex plane. It is the minimal
doubling of $\R$ that makes $-\bone$ a square, and the new element appears
necessarily as the second component of a pair, not as a mysterious primitive
imported from outside.\marginnote{\footnotesize Cayley--Dickson produces
$\R \to \C \to \mathbb{H} \to \mathbb{O}$, doubling dimension at each step and
losing one algebraic property: commutativity, associativity, alternativity. The
process has no known terminal object.}

Write a complex number as a pair $(x, y)$ with $x, y \in \R$, and define
$\Exp(x, y)$ by the same power series, with powers computed via Cayley--Dickson
multiplication. This is a \emph{pronounceable} definition: every term is a
finite unordered collection of pairs, computed by a structural rule requiring no
appeal to geometry or trigonometry.

Evaluating on the imaginary axis --- setting $x = 0$ and letting $y$ vary ---
the series computes:
\[
  \Exp(0, y) \;=\; (\cos y,\; \sin y).
\]
This function is periodic. Its period is the smallest positive $y$ for which
$\Exp(0,y)$ returns to the multiplicative identity of the doubled structure:
\[
  \Exp(0, 2\pi) \;=\; (\bone, \bzero).
\]
One full period returns you to the unit. That is all Euler's identity says ---
and stated this way, it is not an identity at all but a \emph{period theorem}:
the exponential type constructor, evaluated at its own period, returns $\bone$.
The period happens to be $2\pi$, and $2\pi$ is transcendental.

% ======================================================
\section{Transcendence as syntactic incompleteness}

We now have the following situation. The function $\Exp(0, y)$ is defined
entirely within the algebraic/type-theoretic framework: a power series with
rational coefficients, evaluated in a structure obtained by a purely structural
doubling construction. Every finite partial sum is an algebraic expression. The
series converges, the resulting function is periodic, and its period is
witnessed by
\[
  \Exp(0, 2\pi) \;=\; (\bone, \bzero).
\]
But $2\pi$ is \emph{inexpressible} in the language used to define it. No
polynomial with rational (or algebraic) coefficients has $\pi$ as a root. The
object has generated, from within itself, a value that escapes its own
definability.

The standard form $e^{i\pi} = -1$ obscures this completely. It privileges the
half-period, smuggles in $i$ as a geometric primitive, and produces a
relationship between three ``mysterious'' constants that looks like a
coincidence requiring explanation. The form $\Exp(0, 2\pi) = (\bone, \bzero)$
reveals it as a period theorem: the exponential evaluated at its own period
returns the unit. The mystery was entirely ecological.

\medskip\noindent\textbf{This is G\"{o}del's phenomenon, one level down.}

\medskip\noindent
G\"{o}del showed that a sufficiently expressive formal system generates true
statements it cannot prove --- statements whose truth exceeds the system's
ability to certify it from within. Here, a sufficiently expressive type
constructor generates a periodic function whose period exceeds the system's
ability to express it algebraically. In both cases the system is not broken; it
is simply being asked to describe something larger than
itself. Transcendence is the arithmetic shadow of
incompleteness.\marginnote{\footnotesize More precisely: $\pi$ is transcendental
(Lindemann, 1882). The equation $\Exp(0,\pi) = (-\bone, \bzero)$ together with
Lindemann--Weierstrass gives this immediately.}

This also illuminates \emph{why} polynomials with integer or rational
coefficients form a valid ground for transcendence theory. They are exactly the
expressions generated by the algebraic type constructors applied to $\bone$ and
$\btwo$; they are the language of finite algebraic data. Transcendental numbers
are those unreachable from $\bone$ by any finite composition of algebraic
operations --- the values that no algebraic data type can name.

% ======================================================
\section{A remark on bits}

We close with a small observation connecting to information theory without
developing it fully. A \emph{bit} is $\bone/\btwo$: a choice of one from two,
a selection of an inhabitant of $\btwo$, the minimum unit of distinction. The
\emph{information content} of a fair binary choice is
$\log_2(\btwo/\bone) = \log_2 2 = 1$ --- a Shannon bit. These two quantities
are related by the $\log/\Exp$ adjunction, which is the same adjunction
separating the algebraic from the transcendental: $\Exp$ generates the
transcendental from the algebraic, and $\log$ inverts it.

The bit is algebraic. Information --- measured in bits --- is transcendental.
Or rather: the \emph{unit} of information is algebraic ($\bone/\btwo$), but the
\emph{measure} requires the logarithm, the inverse of the exponential type
constructor that generated the transcendental in the first place. There is a
small closed universe here, and its boundary is $\Exp$.

\bigskip
\noindent\textbf{Acknowledgements.} Thanks are due to dialogue itself --- and
in anticipation to Conor McBride, who will no doubt find the sharpest way to
say all of this.

\end{document}